\section{循环群}

由一个元素$g$生成的群称为循环群，记为$G=\langle g\rangle$。其中每个元素都为$g$的整数次幂，且
\[
g^rg^s=g^{r+s}=g^{s+r}=g^sg^r,
\]
所以循环群一定交换。整数加法群$\Z$由1生成，是无限循环群的基本例子。先弄清$\Z$的子群，就可以通过同态研究一般循环群。

\begin{proposition}
$\Z$的每个子群都形如$m\Z$，其中$m$为唯一的非负整数。对$m,n\geq0$，
\[
n\Z\supseteq m\Z\quad\Longleftrightarrow\quad n\mid m.
\]
\end{proposition}
\begin{proof}
设$H\leq\Z$。若$H=\{0\}$，取$m=0$。否则$H$含非零整数；同时含其负元，所以必含正整数。取$H$中最小的正整数$m$。

对任意$x\in H$，作带余除法：
\[
x=qm+r,\qquad q,r\in\Z,\quad0\leq r<m.
\]
因为$x,m\in H$，故$r=x-qm\in H$。若$r>0$，就得到$H$中比$m$更小的正整数，矛盾。因此$r=0$，即$x\in m\Z$。这说明$H\subseteq m\Z$；反过来，$m\in H$使$m$的全部整数倍都属于$H$，故$H=m\Z$。非零子群的$m$由最小正整数唯一确定。

若$n\Z\supseteq m\Z$，则$m\in n\Z$，即$n\mid m$；若$m=nt$，则每个$mk=n(tk)$都在$n\Z$中，得到反向推论。含0的情形也按$m=nt$解释。
\end{proof}

\begin{theorem}
无限循环群同构于$\Z$，其子群也均为循环群。若$G=\langle g\rangle$的阶为$m$，则$G\cong\Z/m\Z$，且对$m$的每个正因子$s$，恰有一个$s$阶子群，为
\[
\langle g^{m/s}\rangle.
\]
\end{theorem}
\begin{proof}
定义满同态
\[
\varphi:\Z\longrightarrow G,\qquad n\longmapsto g^n.
\]
满射性由$g$生成$G$得到；它保持运算，因为$g^{r+s}=g^rg^s$。
其核是$\Z$的子群，所以形如$m\Z$，$m\geq0$。

若$m=0$，$\varphi$单射，故$G\cong\Z$。此时子群也对应于$n\Z$，所以仍为循环群，生成元为$g^n$。

若$m>0$，同态基本定理给出$G\cong\Z/m\Z$，从而$|G|=m$。$m$又是使$g^m=e$的最小正整数，所以它就是$g$的阶。
由子群对应定理，$G$的子群对应于包含$m\Z$的子群$t\Z$，即$t\mid m$。其像为$\langle g^t\rangle$，而
\[
e,g^t,g^{2t},\ldots,g^{(m/t-1)t}
\]
两两不同，且已包含全部$t$的倍数指数所给出的元素。因此其阶为$m/t$。若要求阶为$s$，只能取$t=m/s$，所以所述子群存在且唯一。
\end{proof}

上述分类还可用来计算有限阶元素的幂的阶。

\begin{proposition}
设$g$的阶为正整数$m$。则$g^l=e$当且仅当$m\mid l$，并且对任意整数$s$，
\[
\operatorname{ord}(g^s)=\frac{m}{\gcd(m,s)}.
\]
因此$m$阶循环群的生成元恰为满足$\gcd(m,s)=1$的$g^s$。不同的生成元有$\varphi(m)$个，其中$\varphi(m)$表示$0\leq s<m$中与$m$互素的整数个数。
\end{proposition}
\begin{proof}
同态$\varphi:\Z\to\langle g\rangle$的核为$m\Z$，所以$g^l=e$等价于$l\in m\Z$。

由循环群的子群分类，存在$m$的正因子$d$，使$\langle g^s\rangle=\langle g^d\rangle$。两个生成元互相属于对方生成的子群，因此有整数$n,v,t,w$满足
\[
s=nd+vm,\qquad d=ts+wm.
\]
第一式连同$d\mid m$说明$d\mid s$，所以$d\mid\gcd(m,s)$。第二式说明$\gcd(m,s)\mid d$。故$d=\gcd(m,s)$。
由前一定理，$\langle g^d\rangle$的阶为$m/d$，得到所需阶公式。

最后，$g^s$生成整个$m$阶群当且仅当其阶为$m$，也就是$\gcd(m,s)=1$。指数只须按模$m$区分，故其个数为$\varphi(m)$。
\end{proof}

下面寻求有限交换群成为循环群的判别条件。先利用交换性，将互素阶元素合成一个阶为两者乘积的元素。

\begin{lemma}
若群元素$g,h$交换，其阶分别为互素的$m,n$，则$gh$的阶为$mn$。
\end{lemma}
\begin{proof}
由交换性，$(gh)^{mn}=g^{mn}h^{mn}=e$，所以$gh$有有限阶，记为$l$，则$l\mid mn$。
另一方面，由$(gh)^l=e$再取$m$次幂，得到
\[
e=(gh)^{lm}=g^{lm}h^{lm}=h^{lm}.
\]
因此$n\mid lm$。因为$\gcd(m,n)=1$，有$n\mid l$。同样取$n$次幂可得$m\mid l$，故$mn\mid l$。合并两边整除关系，得到$l=mn$。
\end{proof}

\begin{theorem}\label{p2:max-order}
有限交换群$G$中存在一个元素，其阶为$G$中所有元素的阶的公倍数。
\end{theorem}
\begin{proof}
在有限集合$G$中取阶最大的元素$g$，记其阶为$n$。要证明任一元素的阶均整除$n$。

反设存在$g_1\in G$，其阶$n_1$不整除$n$。由整数的素因子分解，必有某个素数$p$在$n_1$中的指数大于在$n$中的指数。写成
\[
n_1=p^rl,\qquad n=p^sm,\qquad r>s,\quad p\nmid l,\quad p\nmid m.
\]
由幂的阶公式，
\[
\operatorname{ord}(g_1^l)=p^r,\qquad
\operatorname{ord}(g^{p^s})=m.
\]
这两个元素因$G$交换而交换，阶又互素。由引理，
\[
\operatorname{ord}(g_1^l g^{p^s})=p^rm>p^sm=n.
\]
这与$g$的阶最大矛盾，所以不存在这样的$g_1$，结论成立。
\end{proof}

\begin{theorem}\label{p2:cyclic-root-count}
有限交换群$G$为循环群，当且仅当对每个正整数$m$，方程$x^m=e$在$G$中至多有$m$个解。
\end{theorem}
\begin{proof}
\textbf{充分性：}由前一定理，存在阶为$n$的$g\in G$，使全体元素的阶都整除$n$。所以每个$x\in G$都满足$x^n=e$。由题设，$|G|\leq n$。
但$\langle g\rangle$本身已有$n$个元素，且包含于$G$，故$|G|=n$、$G=\langle g\rangle$。

\textbf{必要性（补充证明）：}设$G=\langle g\rangle$的阶为$n$，令$d=\gcd(n,m)$，写$n=dn_0$、$m=dm_0$，其中$\gcd(n_0,m_0)=1$。则
\begin{align*}
(g^k)^m=e
&\iff n\mid km\\
&\iff n_0\mid km_0\\
&\iff n_0\mid k.
\end{align*}
在$0\leq k<n$中，这样的指数恰为$0,n_0,\ldots,(d-1)n_0$，共$d$个。因为$d\leq m$，方程的解数至多为$m$。
\end{proof}
